\documentclass[tikz,border=3mm]{standalone}
% Spolecny styl vlastnich vektorovych ilustraci; preklad jednotlivych stackup_*.tex.
\usepackage[czech]{babel}
\usepackage[utf8]{inputenc}
\usepackage[T1]{fontenc}
\usepackage{tgheros}
\renewcommand{\familydefault}{\sfdefault}
\definecolor{sig}{RGB}{244,185,75}
\definecolor{gnd}{RGB}{107,183,141}
\definecolor{pwr}{RGB}{228,129,112}
\definecolor{frfour}{RGB}{245,237,214}
\tikzset{copper/.style={draw=black!55,line width=0.4pt},
 layerlabel/.style={anchor=east,font=\small},
 note/.style={font=\small,align=center}}

\usetikzlibrary{arrows.meta}
\begin{document}
\begin{tikzpicture}[font=\small,>=Latex]
 \foreach \offset/\variant in {0/0,7/1} {
  \begin{scope}[xshift=\offset cm]
   \draw[red!70!black,line width=4pt] (0,1.5)--(1,1.5);
   \filldraw[fill=pwr,draw=black!60] (.65,.95) rectangle (1.35,2.05);
   \draw[black,line width=1.2pt] (1.35,1.5)--(2.65,1.5);
   \filldraw[fill=black!65,draw=black] (1.5,1.05) rectangle (2.5,1.95);
   % Kratke siroke pripojeni, bez spolecneho uzkeho hrdla.
   \filldraw[fill=gnd,draw=black!55] (2.65,1.05)--(3.35,1.05)--(3.9,.6)
     --(5.6,.6)--(5.6,2.4)--(3.9,2.4)--(3.35,1.95)--(2.65,1.95)--cycle;
   \draw[black!60] (2.65,.95) rectangle (3.35,2.05);
   \node at (2,2.6) {C1};
   \node[font=\scriptsize] at (1,.55) {PWR ploška};
   \node[font=\scriptsize] at (3,.55) {GND ploška};
   \ifnum\variant=0
    \filldraw[fill=gnd!55!black,draw=black!65] (4.75,1.5) circle (.61);
    \fill[white] (4.75,1.5) circle (.36);
    \draw[blue!75!black,->,very thick] (3.45,1.5)--(4.05,1.5);
    \node at (3,-.25) {\textbf{Jeden větší GND prokov}};
   \else
    \foreach \x in {4.15,5.1} {
     \foreach \y in {1,2} {
      \filldraw[fill=gnd!55!black,draw=black!65] (\x,\y) circle (.25);
      \fill[white] (\x,\y) circle (.11);
     }
    }
    \draw[blue!75!black,->,thick] (3.4,1.5)--(3.9,1.05);
    \draw[blue!75!black,->,thick] (3.4,1.5)--(3.9,1.95);
    \draw[blue!75!black,->,thick] (4.4,1.5)--(4.88,1.08);
    \draw[blue!75!black,->,thick] (4.4,1.5)--(4.88,1.92);
    \node at (3,-.25) {\textbf{Více menších GND prokovů}};
   \fi
  \end{scope}
 }
 \node[font=\scriptsize,anchor=west] at (0,-.9) {Pohled shora; všechny GND prokovy vedou do téže souvislé zemní plochy pod kondenzátorem.};
 \node[font=\scriptsize,anchor=west] at (0,-1.25) {Stejná vzdálenost ke GND vrstvě; šipky ilustrují rozdělení proudu. Není v měřítku ani určeno pro výrobu.};
\end{tikzpicture}
\end{document}
